\begin{itemize}
\item \textbf{Scale.} One modulus (23), one one-layer architecture, one learning rate ($3\times10^{-3}$, chosen on uniform sampling with an unreported pilot seed, which favours uniform), 4000 steps. A different modulus, learning rate, depth or step budget could change the ordering of the systems.
\item \textbf{Censoring and power.} Most grid cells have all runs censored at 4000 steps, so they carry no information about relative speed. The main cell has 5 seeds and the grid 3; the test registered for H1 is not significant (Table~\ref{tab:tests}), so the data neither establish nor exclude a difference. Means of censored values are lower bounds; we reported the fraction reached alongside.
\item \textbf{Grokking only partly reached.} Within the budget, delayed generalisation was seen at the main cell for uniform sampling and in the cells with $\lambda\ge1$ and $f\ge0.5$; at $\lambda=0$ and at $f=0.4$ it was not reached by any run, and we cannot say if longer runs would.
\item \textbf{Training instability.} Held-out accuracy curves show sharp transient drops (Fig.~\ref{fig:curves}), and one uniform run at the main cell ended with a training accuracy below 0.9. A smaller learning rate with longer training would be a cleaner test than ours.
\item \textbf{Narrow curriculum family.} A single difficulty measure ($\max(a,b)$), a single linear pacing function and rank-based pools; other measures (e.g.\ sum size or residue structure) and exposure-matched baselines were not tried. The explanation for the pacing result is a hypothesis, not tested by an isolating run.
\item \textbf{Reimplementation.} The uniform baseline and the model are our own code; no result was compared to published numbers, and no mechanistic analysis of the learned solution was done.
\item \textbf{Held-out definition.} Held-out accuracy is on the complement of a random split; both ordering systems see exactly the same training pairs as uniform, only their timing differs.
\end{itemize}
